\chapter{過去形 — El Pasado}
\unidadheader{11}{過去形}{El Pasado}{Contar lo que hiciste}

\begin{tcolorbox}[vocabbox, title={📌 Vocabulario Nuevo}]
\begin{tabularx}{\linewidth}{llX}
\toprule
\textbf{Japonés} & \textbf{Español} & \textbf{Tipo} \\
\midrule
\vocabitem{見る}{みる}{ver / mirar}{v. gr.2}
\vocabitem{聞く}{きく}{escuchar / preguntar}{v. gr.1}
\vocabitem{買う}{かう}{comprar}{v. gr.1}
\vocabitem{読む}{よむ}{leer}{v. gr.1}
\vocabitem{書く}{かく}{escribir}{v. gr.1}
\vocabitem{話す}{はなす}{hablar / contar}{v. gr.1}
\vocabitem{泳ぐ}{およぐ}{nadar}{v. gr.1}
\vocabitem{走る}{はしる}{correr}{v. gr.1}
\vocabitem{遊ぶ}{あそぶ}{jugar / divertirse}{v. gr.1}
\vocabitem{昨日}{きのう}{ayer}{adv.}
\bottomrule
\end{tabularx}
\end{tcolorbox}

\subsection*{🗣️ Frases Clave}

\frase{きのう\ruby{何}{なに}をしましたか？}{¿Qué hiciste ayer?}
\frase{\ruby{映画}{えいが}を\ruby{見}{み}ました。}{Vi una película.}
\frase{あまり\ruby{楽}{たの}しくなかったです。}{No fue muy divertido.}
\frase{\ruby{先週}{せんしゅう}どこへ\ruby{行}{い}きましたか？}{¿A dónde fuiste la semana pasada?}

\subsection*{⚙️ Gramática}

\patron{Pasado afirmativo y negativo}
  {〜ました / 〜ませんでした}
  {\ej{\ruby{図書館}{としょかん}で\ruby{本}{ほん}を\ruby{読}{よ}みました。}
    {Leí un libro en la biblioteca.}
  \ej{きのう\ruby{学校}{がっこう}に\ruby{行}{い}きませんでした。}
    {Ayer no fui a la escuela.}}

\patron{Pregunta en pasado}
  {〜ましたか？}
  {\ej{あの\ruby{映画}{えいが}を\ruby{見}{み}ましたか？}
    {¿Viste esa película?}
  \ej{はい、\ruby{見}{み}ました。/ いいえ、\ruby{見}{み}ませんでした。}
    {Sí, la vi. / No, no la vi.}}

\begin{tcolorbox}[ejerciciobox, title={📝 Ejercicios Rápidos}]
\begin{enumerate}
  \item Conjuga en pasado negativo: \ruby{書}{か}く → ?
  \item ¿Cómo preguntas qué hizo alguien el fin de semana?
  \item Describe qué hiciste ayer en 2 oraciones.
\end{enumerate}
\vspace{0.4em}
\textbf{Respuestas:}
1) \ruby{書}{か}きませんでした\quad
2) \ruby{週末}{しゅうまつ}に\ruby{何}{なに}をしましたか？\quad
3) (personal)
\end{tcolorbox}

\begin{tcolorbox}[kanjibox, title={🀄 Kanjis de la Unidad}]
\begin{tabularx}{\linewidth}{cllXX}
\toprule
\textbf{Kanji} & \textbf{On} & \textbf{Kun} & \textbf{Significado} & \textbf{Vocab} \\
\midrule
\kanjirow{見}{けん}{み}{ver / mostrar}{\ruby{見}{み}る / \ruby{見物}{けんぶつ}}
\kanjirow{聞}{ぶん・もん}{き}{escuchar / preguntar}{\ruby{聞}{き}く / \ruby{新聞}{しんぶん}}
\kanjirow{買}{ばい}{か}{comprar}{\ruby{買}{か}う / \ruby{買い物}{かいもの}}
\kanjirow{読}{どく}{よ}{leer}{\ruby{読}{よ}む / \ruby{読書}{どくしょ}}
\kanjirow{書}{しょ}{か}{escribir}{\ruby{書}{か}く / \ruby{教科書}{きょうかしょ}}
\bottomrule
\end{tabularx}
\end{tcolorbox}
